% Pursuing Stacks (cotes 134-2 to 134-8): archive cote 134-4, site page 45 (PDF page 46), typescript page 133.
% typescript. Draft compared: ps3.tex:2271-2344.
% AI-transcribed by Claude Opus 5.5 (claude-opus-5-5, adaptive thinking, effort high), Message Batches API, batch ps2 (msgbatch_01NVwr1Yuo8qJf5rykeGBXPK), 2026-11-12; not yet checked by eye.
\scanpage{134-4/45}
\mspage{133}

which is not always the case. (Take for instance M to be an abelian cate-%
gory,  say projective complexes of modules and quasi-isomorphism between
these;  in this case,  more generally whenever M has a ``zero object'' name-%
ly one which is both initial and final,  any interval in M is trivial i.e.
$\hand{\delta}_0=\hand{\delta}_1$ ,  and hence \xout{the correspondin} any weak homotopy interval structure
on M is trivial namely $W_h$ is reduced to isomorphisms ...).

Thus the more reasonable question here is whether any homotopy struc-%
ture h on M,  definable in terms of a weak homotopy interval structure,
\inserted[typed]{and such that $W_h\hand{\subset}W$,} \inserted{Clearly,}
satisfies $h\hand{\leq}h_W$ . \over{F}{f}or such an h,  any weak homotopy interval \underline{I} (for  h)
\deleted{clearly} satisfies  $u(\hand{\delta}_0)=u(\hand{\delta}_1)$ ,  where u: $M\hand{\rightarrow}W^{-1}M$  is the canonical
functor  (indeed,  it is enough \xout{that} for this that \deleted{\xout{W\illegible{1}}} $W_h\hand{\subset}\hand{\overline{W}}$  instead
of $W_h\hand{\subset}W$,  where $\hand{\overline{W}}$ is the strong saturation of W), \inserted{and conversely if $W=\overline{W}$} \deleted{\xout{To deduce that}}
\inserted{and \deleted{conversely} if moreover u \uline{commutes to finite products}.}
On the other hand,  $h\hand{\leq}h_W$ means that \xout{xxx} any $\underline{I}\hand{\in}\hand{\Sigma}_h$ is in $\hand{\Sigma}_{(h_W)}$ , which \margin{$\times$}
also means that it's endpoint\xout{x} \xout{xxx} sections $\hand{\delta}_0$ , $\hand{\delta}_1$ are  $h_W$ -homotopic,
namely may be joined by a chain of sections,  any two consecutive of
which are related by \deleted{a}\inserted{some} \underline{J}-homotopy,  where \underline{J} is a W-\uline{aspheric interval}.
Thus we get the:

\uline{Proposition}  Let W a saturated set of arrows in M (M stable under finite
products),  hence a corresponding homotopy structure $h_W$ on M,  defined in
terms of \inserted{W-}aspheric intervals\hand{*}in M as a generating set of weak homotopy \margin{* 11.4. Cf. note on preceding page.}
intervals for $h_W$ .  \xout{The} Consider the following conditions:

(i)  \inserted{$h_W$} \deleted{\xout{xxxxxxx}}is the widest of all homotopy structures \xout{xx} h on M which are
(a)  \xout{xxmpatible} strictly compatible with W  i.e. such that $W_h\hand{\subset}W$  and \inserted{moreover}
(b)  definable in terms of a weak homotopy interval structure.

(ii)  For any object I of M and two sections $\hand{\delta}_0,\hand{\delta}_1$ of I such that
$u(\hand{\delta}_0)=u(\hand{\delta}_1)$  (where u: $M\hand{\rightarrow}W^{-1}M$ is the canonical functor), $\hand{\delta}_0$ and $\hand{\delta}_1$ are
$h_W$-homotopic,  namely can be joined by a chain of elementary homotopies
as above,  involving \xout{asp} W-\uline{aspheric}\hand{*} intervals \underline{J} .

Then (ii) implies (i),  and conversely if W is strongly saturated\hand{\unsure{,}}
and moreover u commutes to finite products.

In connection with the $\hand{\pi}_0$-functor, \xout{xxxwhich} we are going to get pretty
natural conditions \xout{xxsuring} in terms of 0-connectedness for insuring (ii)
and hence (i),  which should apply I guess to all ``reasonable'' modelizers
(M,W).  \xout{In practical terms,}  It would thus seem that  in practical terms,
the definition of $h_W$ \inserted{is} the best,  in \hand{all} cases of actual interest to us.
Of course,  in case W is strongly saturated and u commutes with finite
products,  the widest h is the one \xout{defined by} whose weak homotopy inter-%
vals are the triples $(I,\hand{\delta}_0,\hand{\delta}_1)$  satisfying $u(\hand{\delta}_0) = u(\hand{\delta}_1)$,  which \over{x}{w}e could
have used instead of just W-aspheric intervals,  \xout{The advantage} which are

\margin{45}
